This report asked whether a patch-based forecaster preserves the frequency content of signals that classical sampling theory says are fully representable, and how $f_s$, $P$ and $S$ decide when it does not. The characterisation is complete and depends on no fit. A tone completing a whole number of cycles in one stride repeats its tokens; one completing a whole number in one patch repeats their content; and the union of the two combs is a closed-form candidate set computable for a geometry never measured. That set locates where a loss could sit, not that one exists, and the report holds the two apart throughout.

Whether a trained model does lose information there remains open. The design that decides it is fixed in advance, and it returns not reportable or not identified rather than a null when a fit fails to earn a verdict. Two consequences already follow. Patch geometry is a deployment parameter to be chosen against the known spectral content of a plant, not a default to inherit; and since a detector placed after tokenisation inherits the blind spot it is meant to find, monitoring must act on the raw signal upstream. Neither waits on the fits, and neither is available to a system that treats the tokeniser as opaque.

Three extensions would sharpen the claim: independently trained instances per geometry, which a single seed cannot separate from geometry; an untrained-network baseline on the same dispersion measurement; and a test of whether a candidate set derived at $512$\,Hz survives translation to telemetry sampled once a minute.
